\section{Exact outputs and finite arithmetic}\label{app:output-separation}
Three different assertions should not be conflated. The existence of the
relative real atlas is a geometric theorem. The first-order procedures
produce spectral or entrance outputs from algebraic presentations.
The rational calculations of specific Fisher programs are finite
specializations, not implementations or proofs of the entire atlas.

For the direct procedures the input is a finite list of polynomials with
rational or real algebraic coefficients, a Boolean formula of polynomial
sign conditions and the quantifiers specifying the model and its maps.
A real algebraic coefficient is represented by its integer polynomial
and a rational isolating interval. Rational powers use a nonnegative
auxiliary variable and its integer-power equation. Compact minimum
formulas have the form
\[
 \exists x\,[\Phi(x)\wedge Q(x)=z]
 \ \wedge\ \forall y\,[\Phi(y)\Rightarrow Q(y)\ge z].
\]
The sign or limit formulas retain the centre as a free fixed parameter.
Elimination produces a finite output formula; at algebraic parameter
values, root isolation gives the numerical algebraic outputs. In
particular, the quadratic-probe graph and the entrance minimum graph fit
this encoding without any projective blow-up representation.

The separate certificate search in Appendix~\ref{app:effective} concerns
geometric presentations. Its existence-by-enumeration role is
supplementary: the direct outputs above do not wait for that search,
and no run of the finite example regression is described as having
constructed such a certificate. This separation leaves the geometric
proofs and the direct algorithmic conclusions intact.

For the two cubic critical examples, use the coordinate orders
\[
 (a_0,b_0,X,c_0,w,k,u_f,u_h,v_f,v_h)
 \quad\hbox{and}\quad
 (a_0,b_0,X,c_0,d_0,Y,k,u_f,u_h,v_f,v_h),
\]
respectively. At the endpoint, $w=-(d_1+d_2)\ge0$; for a first-order
representative choose $d_1=d_2=-w/2$. The nonnegative coordinates are
$(a_0,X,w)$ and $(a_0,X,Y,k)$. Both displayed specializations have the
empty cone support, so $w=0$ at their endpoint optimum and the endpoint
split is uniquely zero. Thus the representative choice introduces no
second-order ambiguity into those exact examples.

Let $L$ and $D$ be respectively the seven free score columns and the cone
columns, $W=\diag(w_j/P_{je})$, and let $B$ be the base probability
velocity. All entries are rational at the stated examples. Set
\[
 G_L=L^{\mathsf T}WL,\quad
 b_f=G_L^{-1}L^{\mathsf T}WB,\quad
 \widetilde D=D-LG_L^{-1}L^{\mathsf T}WD,
\]
\[
 Q=\widetilde D^{\mathsf T}W\widetilde D,\quad
 g=\widetilde D^{\mathsf T}W(B-Lb_f),\quad
 b_*=(B-Lb_f)^{\mathsf T}W(B-Lb_f).
\]
The exact checks give $\rank L=7$, $\rank Q=2$, $g<0$ entrywise and
nonnegative principal minors of $Q$. Every nonempty independent-support
KKT test fails, whereas the empty support succeeds. Moreover
$z^{\mathsf T}Qz-2g^{\mathsf T}z>0$ for every nonzero $z\ge0$.
This proves uniqueness of the cone optimum despite the singular
projected Gram matrix. The free optimum is $b_f$ and $\mu^2=b_*/4$.
Multiplication and division of the exact probability jets then give
$\mu k$ by \eqref{eq:exact-second-coefficient}.
For positive rational interval endpoints $l,u$, the tests for $\mu$ and
$k$ are respectively
\[
 l^2<\mu^2<u^2,
 \qquad l^2\mu^2<(\mu k)^2<u^2\mu^2,
 \quad\mu k>0.
\]
After clearing positive denominators these are integer comparisons.
The same exact arithmetic forms and inverts the two Fisher matrices in
\eqref{eq:M-quarter}--\eqref{eq:M-eighth}. The ancillary record gives
the full rational Gram matrices, linear vectors, optimizers and interval
tests; reproducibility of these finite operations is separate from the
mathematical arguments above.
